The linear multiplicity-one problem reduces by Frobenius descent to the quiver
action of Section~\ref{sec:quiver}. The regular rank-zero cases
are handled by the two-leg theorem, while the scalar rank-zero
case requires a separate identity descendant. We prove these statements before treating all nine rational families
and completing the linear multiplicity-one argument.

\subsection{The residual action}\label{m1:subsec:reduction-to-GL4}
For $G_{\mathrm L}=\GL_4(F)\times\GL_2(F)$ and the subgroup $H_{\mathrm L}$
of Theorem~\ref{m1:thm:main}, the group $J=H_{\mathrm L}\times H_{\mathrm L}$ acts by
\[
 ((a_1,b_1),(a_2,b_2))(g,z)
 =\bigl(\diag(a_1,b_1)g\diag(a_2,b_2)^{-1},\ b_1zb_2^{-1}\bigr).
\]
Componentwise transpose preserves $H_{\mathrm L}$ and normalizes this action, sending
the two acting factors to their inverse transposes in reversed order.
The projection $(g,z)\mapsto z$ is equivariant with transitive base
$\GL_2(F)$. At $z=I_2$ its stabilizer is
\[
 J_{I_2}=\{((p,q),(u,q)):p,q,u\in\GL_2(F)\}
 \simeq L(F),\qquad L=\GL(P)\times\GL(Q)\times\GL(U),
\]
where $P,Q,U$ have dimension two. Both $J$ and $J_{I_2}$ are unimodular.
The fiber is $\cX\simeq\GL_4(F)$ and transpose preserves its base point.
The finite extensions are unimodular as well: a modulus is trivial on
the ordinary group and factors through a finite group. Thus
\begin{equation}\label{m1:eq:residual-Frobenius}
 \mathcal D(G_{\mathrm L})^J\simeq\mathcal D(\cX)^{L(F)}
\end{equation}
is transpose-compatible, including its odd component. In block
coordinates $x=(A,B,C,D)$, the resulting action and involution are
exactly \eqref{m1:eq:residual-action}--\eqref{m1:eq:tau}.
The variety $\cX$ is smooth affine and $\alpha$ is algebraic over $F$.
It remains to prove $\mathcal D(\cX)^{\widetilde L,\varepsilon}=0$.
By Corollary~\ref{m1:cor:HC-invariance}, this follows from rational transpose
stability of every closed orbit and the absence of sign-equivariant
distributions on its normal space.
Only one representative per rational orbit is needed: $y=l x$ identifies
both stabilizers and normal spaces by conjugation and the differential of
$l$. The geometric data were established in Section~\ref{sec:quiver};
it remains to rule out sign-equivariant distributions on these normal spaces.

\subsection{The two-leg distribution theorem}\label{m1:sec:two-leg}
\subsubsection{Adjoints and the statement}\label{m1:subsec:two-leg-statement}

Let $V$ and $W$ be two-dimensional $F$-vector spaces, and let
$T\subset \GL(W)$ be a maximal torus.  Put \(\Gamma=\GL(V)\times T\) and \(\Vlegs=\Hom(W,V)\oplus
\Hom(V,W)\). We write a point of $\Vlegs$ as $(X,Y)$, with \(X:W\longrightarrow V\) and
\(Y:V\longrightarrow W\), and let $\Gamma$ act by
\[
    (a,t)\cdot(X,Y)=\bigl(aXt^{-1},\,tYa^{-1}\bigr).
\]
Fix a basis of $V$ with its standard symmetric form. Choose a
nondegenerate symmetric matrix $S$ on $W$ such that
\begin{equation}\label{m1:eq:torus-self-adjoint}
 S t=t^t S\qquad(t\in T).
\end{equation}
Define the adjoint exchange
\begin{equation}\label{m1:eq:two-leg-adjoint}
 \tau=\tau_S,\qquad \tau_S(X,Y)=(Y^tS,S^{-1}X^t).
\end{equation}
It satisfies $\tau_S^2=1$ and, by direct multiplication,
\[
 \tau_S\bigl((a,t)\cdot(X,Y)\bigr)
 =(a^{-t},t^{-1})\cdot\tau_S(X,Y).
\]
In particular it normalizes the given action. Such a form exists over $F$:
in a diagonal realization of a split torus take $S=I_2$; for a nonsplit
torus identify $W$ with a quadratic field $K$ and $T$ with $K^\times$
acting by multiplication, and take the Gram matrix of
$\operatorname{Tr}_{K/F}(ww')$. The trace form is symmetric and
nondegenerate because $K/F$ is separable, and multiplication is self-adjoint.

A change of basis $w=Cw'$ replaces $(X,Y,S)$ by
$(XC,C^{-1}Y,C^tSC)$ and conjugates $T$ accordingly, preserving \eqref{m1:eq:torus-self-adjoint} and
\eqref{m1:eq:two-leg-adjoint}. In a fixed realization, any other such form is
$S=S_0c$ with $c\in T$: indeed $S_0^{-1}S$ commutes with $T$, whose
centralizer in $\operatorname{End}_F(W)$ is its quadratic \'{e}tale algebra.
Then $\tau_S=(1,c^{-1})\tau_{S_0}$, so these forms induce the same
operator on $\Gamma$-invariant distributions. It therefore suffices to prove
the split and trace-form realizations just specified. In this section,
the adjoint exchange on the two legs always means $\tau_S$; on
$\operatorname{End}_F(W)$ it means $A^*=S^{-1}A^tS$.

Let \(\widetilde \Gamma=\Gamma\rtimes\langle\tau\rangle\) and let $\varepsilon$ be the sign character,
\(\varepsilon|_\Gamma=1\) and \(\varepsilon(\tau)=-1\).

\begin{theorem}[Two-Leg Lemma]\label{m1:thm:two-leg}
For every maximal torus $T\subset \GL_2(F)$ and every nondegenerate symmetric form $S$ satisfying
\eqref{m1:eq:torus-self-adjoint},
\[
    \mathcal D(\Vlegs)^{\widetilde \Gamma,\varepsilon}=0.
\]
Equivalently, \(\mathcal D(\Vlegs)^\Gamma=\mathcal D(\Vlegs)^{\Gamma,\tau}\). \end{theorem}

We first establish the matrix-axis boundary statement and
the rank-one factorization result common to both rational torus types.

\subsubsection{Matrix axes and local zeta distributions}\label{m1:subsec:matrix-axes}

Let
\[
    B=
    \bigl(\Hom(W,V)\oplus\{0\}\bigr)
    \cup
    \bigl(\{0\}\oplus\Hom(V,W)\bigr).
\]
Thus $B$ is the union of the two matrix axes.

\begin{lemma}[Matrix two-axis boundary]\label{m1:lem:matrix-boundary}
For every maximal torus $T\subset \GL(W)$,
\[
    \mathcal D(B)^\Gamma=\mathcal D(B)^{\Gamma,\tau}.
\]
Moreover, if \(\lambda(a,t)=\left|\frac{\det a}{\det t}\right|\), then the invariant distributions on the
two rank-two open axes admit
finite-part extensions $P_X$ and $P_Y$ satisfying
\begin{equation}\label{m1:eq:matrix-finite-parts}
    \tau P_X=P_Y,
    \qquad
    gP_X-P_X=-\log\lambda(g)\,R,
    \qquad
    gP_Y-P_Y=+\log\lambda(g)\,R,
\end{equation}
where $R=c\delta_0$ with $c\neq0$.
\end{lemma}

\begin{proof}
\emph{Rank-one axes.} On the rank-one stratum of the $X$-axis, the map \(X\longmapsto \operatorname{im}X\in
\mathbf P(V)\) is $\Gamma$-equivariant, and $\GL(V)$ acts transitively on $\mathbf P(V)$.
At the line $Fe_1$, Frobenius descent leaves the group
\[
    P\times T,
    \qquad
    P=
    \left\{
    \begin{pmatrix}
        \alpha&b\\0&\beta
    \end{pmatrix}:
    \alpha,\beta\in F^\times,\ b\in F
    \right\},
\]
acting on $W^\vee\setminus\{0\}$ by \(c\longmapsto \alpha c t^{-1}\). With our convention for
modular characters,
\[
    \Delta_P
    \begin{pmatrix}
        \alpha&b\\0&\beta
    \end{pmatrix}
    =\left|\frac{\alpha}{\beta}\right|.
\]
Since $\Gamma$ is unimodular, Proposition~\ref{m1:prop:Frobenius-descent} identifies $\Gamma$-invariant
 distributions on the rank-one stratum with distributions on the descended
fiber transforming by the character $\Delta_P^{-1}=|\beta/\alpha|$.  Now take
\[
    p_\beta=\begin{pmatrix}1&0\\0&\beta\end{pmatrix},
    \qquad t=1.
\]
This element acts trivially on the descended geometric fiber, whereas the
required equivariance character is $|\beta|$.  Choosing
$\beta\in F^\times$ with $|\beta|\neq1$ forces the descended distribution to
vanish.  Hence the rank-one $X$-stratum carries no invariant distribution of
the required Frobenius type.  The same argument, after transpose, treats the
rank-one $Y$-stratum.

\emph{Determinant zeta distributions.}\label{m1:para:det-zeta}
We normalize additive Haar measure by
$\operatorname{vol}(M_2(\mathcal O_F))=1$ and, initially for
$\Re(s)\gg0$, define
\[
 Z_X(s)(f)=\int_{M_2(F)}f(X,0)|\det X|^s\,dX,
 \qquad Z_Y(s)=\tau_S Z_X(s).
\]
With the same coordinate Haar measure on the two matrix spaces,
\[
 Z_Y(s)(f)=|\det S|^{s+2}
       \int_{M_2(F)}f(0,Y)|\det Y|^s\,dY,
\]
since $Y=S^{-1}X^t$ has absolute Jacobian $|\det S|^{-2}$.

Theorem~5.13 and Section~5.5 of \cite{HongSun2016}, applied to additive
Haar measure on $\{\det\ne0\}$, give meromorphic continuation and
distribution-valued Laurent coefficients. More precisely, with
$u=s+2$, the family $Z_X(u-2)$ is the case $n=2$, $\chi=1$ of the
matrix zeta integral in \cite[Section~7]{HongSun2016}.
Proposition~7.2(a), together with equation~(30) in its proof, shows that
the pole at $u=0$ is simple and its leading coefficient is a nonzero
multiple of $\delta_0$. Here the Laurent coordinate used in that paper
is $1-q_F^{-u}$, which has derivative $\log q_F\ne0$ at $u=0$.
Thus
\[
 R:=\operatorname*{Res}_{s=-2}Z_X(s)=c\delta_0,\qquad c\ne0.
\]
For the chosen normalization, the lattice formula in the proof of
\cite[Proposition~7.2]{HongSun2016} gives
\[
 I(s):=\int_{M_2(\mathcal O_F)}|\det X|^s\,dX
 =\frac{(1-q_F^{-1})(1-q_F^{-2})}
        {(1-q_F^{-s-1})(1-q_F^{-s-2})},
 \qquad
 c=\frac{(1-q_F^{-1})(1-q_F^{-2})}
        {(1-q_F)\log q_F}.
\]
Because $\tau_S$ fixes the origin, $Z_Y(s)$ has the same residue $R$.

For $g=(a,t)\in\Gamma$, the map $X\mapsto aXt^{-1}$ multiplies
$|\det X|$ by $\lambda(g)$ and has absolute Jacobian $\lambda(g)^2$.
Using the pushforward convention \eqref{m1:eq:distribution-conventions}
and the fact that conjugation by $\tau_S$ inverts $\lambda$, we obtain
\[
 gZ_X(s)=\lambda(g)^{-(s+2)}Z_X(s),\qquad
 gZ_Y(s)=\lambda(g)^{s+2}Z_Y(s).
\]
Define the finite parts using the coordinate $s+2$:
\[
 Z_X(s)=\frac{R}{s+2}+P_X+O(s+2),\qquad
 Z_Y(s)=\frac{R}{s+2}+P_Y+O(s+2).
\]
Then $P_Y=\tau_S P_X$, and comparison of constant terms gives
\[
 gP_X-P_X=-\log\lambda(g)\,R,\qquad
 gP_Y-P_Y=+\log\lambda(g)\,R,
\]
which is \eqref{m1:eq:matrix-finite-parts}.
On each rank-two open axis these finite parts restrict to the nonzero
invariant measure obtained by setting $s=-2$ in the defining integral.


\emph{Finite-part comparison.} Let $D\in\mathcal D(B)^\Gamma$, and let $c_X,c_Y$ be the coefficients of its
restrictions to the two rank-two open axes.  Put \(Q=D-c_XP_X-c_YP_Y\). Then $Q$ is supported on the
union of the two rank-at-most-one axes and
\[
    gQ-Q=(c_X-c_Y)\log\lambda(g)\,R.
\]
The right-hand side is supported at the origin.  Hence the restriction of
$Q$ to either rank-one stratum is invariant, and the rank-one
calculation above forces these restrictions to vanish.  Thus $Q$ is supported
at the origin.  A distribution supported at one point on an $l$-space is a
scalar multiple of the Dirac distribution, so $gQ-Q=0$.  Choosing $g$ with
$\lambda(g)\neq1$ gives $c_X=c_Y$.  Since $\tau P_X=P_Y$ and $\delta_0$ is
fixed by transpose, $D$ is fixed by $\tau$.
\end{proof}

\subsubsection{Singular rank-one factorizations}\label{m1:subsec:rankone-factorization}

\begin{lemma}[Rank-one factorization]\label{m1:lem:rank-one-factorization}
Let $A:W\to W$ have rank one. Suppose that $R$ is an invertible
symmetric matrix satisfying $RA=A^tR$. Put
\[
 \mathcal F_A=\{(X,Y)\in \Vlegs:YX=A\},\qquad
 \tau_R(X,Y)=(Y^tR,R^{-1}X^t).
\]
Then every $\GL(V)\times F^\times$-invariant distribution on
$\mathcal F_A$, with the second factor acting as scalar matrices on $W$,
is invariant under $\tau_R$.
\end{lemma}

\begin{proof}
The relation $RA=A^tR$ makes $\tau_R$ preserve $\mathcal F_A$;
$R^t=R$ makes it involutive. Set $Y'=RY$. Then $Y'X=RA$ is symmetric
of rank one, while the group action is still
$(X,Y')\mapsto(aXt^{-1},tY'a^{-1})$ because $t$ is scalar.
The involution is now $(X,Y')\mapsto((Y')^t,X^t)$.
Choose $C\in\GL_2(F)$ such that
\[
 C^tRA C=\begin{pmatrix}\lambda&0\\0&0\end{pmatrix},\qquad\lambda\ne0.
\]
This is congruence diagonalization of a symmetric rank-one form in
characteristic different from two. The paired changes
$(X,Y')\mapsto(XC,C^tY')$ preserve both the scalar two-leg action and
the ordinary transpose exchange. Relabel these new legs $(X,Y)$.
This congruence argument applies also to nilpotent $A$.

Write
\[
    X=(x_1,x_2),
    \qquad
    Y=\begin{pmatrix}y_1\\y_2\end{pmatrix},
\]
with $x_i\in V$ and $y_i\in V^\vee$. In these paired coordinates, the
defining equation of $\mathcal F_A$ is $YX=\operatorname{diag}(\lambda,0)$,
equivalently
\[
    y_1(x_1)=\lambda,
    \qquad
    y_1(x_2)=0,
    \qquad
    y_2(x_1)=0,
    \qquad
    y_2(x_2)=0.
\]
Consider
\[
    p:\mathcal F_A\longrightarrow
    \Omega_\lambda,
    \qquad
    p(X,Y)=(x_1,y_1),
\]
where \(\Omega_\lambda= \{(v,\ell)\in V\times V^\vee:\ell(v)=\lambda\}\). The group $\GL(V)\times
F^\times$ acts transitively on $\Omega_\lambda$.
Frobenius descent at $(e_1,\lambda e_1^\vee)$ leaves the stabilizer
\[
    H_A=
    \left\{
    \left(
    \begin{pmatrix}t&0\\0&c\end{pmatrix},
    tI_W
    \right):
    t,c\in F^\times
    \right\}
    \simeq(F^\times)^2.
\]
The fiber over $(e_1,\lambda e_1^\vee)$ is
\[
    x_2=xe_2,
    \qquad
    y_2=ye_2^\vee,
    \qquad
    xy=0.
\]
Thus the descended fiber is the scalar two-axis node \(N=\{(x,y)\in F^2:xy=0\}\). The effective quotient
of $H_A$ acts by \(q\cdot(x,y)=(qx,q^{-1}y)\), \(q\in F^\times\). All groups involved are
unimodular. Ordinary transpose sends the base
$(e_1,\lambda e_1^\vee)$ to $(\lambda e_1,e_1^\vee)$. Correct it by
$(a_0,1)$ with $a_0=\operatorname{diag}(\lambda^{-1},1)$. This fixes the
base and sends the fiber coordinates $(x,y)$ to $(y,x)$.
The correcting element composed with transpose is an involution and has
trivial modular character. Prasad's Two-Axis Lemma therefore gives the required
transpose invariance on $N$, and Frobenius descent gives the assertion on
$\mathcal F_A$.
\end{proof}

\subsubsection{The zero-product projective descent}\label{m1:subsec:zero-product-descent}
\begin{lemma}\label{m1:lem:common-U11}
Let $\mathcal F_0=\{(X,Y):YX=0\}$ and
$U_{11}=\{(X,Y)\in\mathcal F_0:\rk X=\rk Y=1\}$.
The decomposition $\mathcal F_0=U_{11}\sqcup B$ has $U_{11}$ open and
$B$ the closed matrix axes. Descent along
$(X,Y)\mapsto\operatorname{im}X=\ker Y\in\mathbf P(V)$ identifies the
ordinary invariant-distribution problem on $U_{11}$ with the
$\eta$-equivariant problem for $P\times T$ on
$(W^\vee\setminus\{0\})\times(W\setminus\{0\})$, where
\begin{equation}\label{m1:eq:common-U11-action}
 a=\begin{pmatrix}\alpha&b\\0&\beta\end{pmatrix},\qquad
 (c,d)\longmapsto(\alpha ct^{-1},\beta^{-1}td),\qquad
 \eta(a,t)=|\beta/\alpha|.
\end{equation}
The corrected odd operator is
\begin{equation}\label{m1:eq:common-U11-odd}
 (c,d)\longmapsto(d^tS,S^{-1}c^t).
\end{equation}
It squares to one, preserves the character $\eta$, and has relative
modular character one on the odd component.
\end{lemma}
\begin{proof}
When both legs have rank one, $YX=0$ forces $\operatorname{im}X=\ker Y$.
The map to $\mathbf P(V)$ is equivariant, and at $Fe_1$ the fiber
is $X=e_1c$, $Y=de_2^\vee$. Matrix multiplication gives
\eqref{m1:eq:common-U11-action}; the upper unipotent entry acts trivially.
The ordinary relative modulus is $\Delta_P^{-1}=|\beta/\alpha|$.
Writing $j=\left(\begin{smallmatrix}0&1\\1&0\end{smallmatrix}\right)$
on $V$, the operator $(j,1)\tau_S$ fixes the base line and gives
\eqref{m1:eq:common-U11-odd}. It has square one and sends the diagonal
parameters $(\alpha,\beta)$ to $(\beta^{-1},\alpha^{-1})$, preserving
$\eta$. The positive relative modular character is one on this
order-two operator, so Proposition~\ref{m1:prop:odd-Frobenius} gives the required invariance.
\end{proof}

\subsubsection{Nonsplit fibers}\label{m1:subsec:two-leg-nonsplit}

Assume that $W=K$, where $K/F$ is a quadratic field extension. We identify \(T=K^\times\subset
\GL_F(K)\) through multiplication and take $S$ to be the trace-form Gram matrix above.
Write $\bar w$ for the nontrivial
Galois conjugate of $w\in K$, and let $\iota(w)=\bar w$.

Every $F$-linear endomorphism of $K$ has a unique expression
\[
    A=m_c+m_d\iota,
    \qquad c,d\in K,
\]
where $m_c$ denotes multiplication by $c$.  For $A=YX$, define
\[
    \pi_{\mathrm{ns}}:\Vlegs\longrightarrow K\oplus F,
    \qquad
    \pi_{\mathrm{ns}}(X,Y)=\bigl(c,N_{K/F}(d)\bigr).
\]
Conjugation by $t\in K^\times$ gives
\[
    t(m_c+m_d\iota)t^{-1}
    =m_c+m_{(t/\bar t)d}\iota,
\]
so $\pi_{\mathrm{ns}}$ is $\Gamma$-invariant.

With respect to the trace pairing \(\langle x,y\rangle_K=\operatorname{Tr}_{K/F}(xy)\), the adjoint of
$A$ is \(A^\ast=m_c+m_{\bar d}\iota\). Hence $\pi_{\mathrm{ns}}$ is fixed by $\tau_S$. Values outside the
norm image have empty fiber. If $N(d)\neq0$, Hilbert 90
shows that the set of $d'$ with $N(d')=N(d)$ is one rational $T$-orbit, with
stabilizer $F^\times$.  If $N(d)=0$, then $d=0$, so the base fiber is a single
point with stabilizer $K^\times$.  These stabilizers are unimodular.

Fix $A=m_c+m_d\iota$ in one such base orbit and put \(T_A=Z_T(A)\). The operators $A$ and $A^\ast$ lie in
the same rational $T$-orbit.  If
$d\neq0$, choose $t_A\in K^\times$ such that \(\frac{t_A}{\bar t_A}=\frac{d}{\bar d}\); if $d=0$, take
$t_A=1$.  Then \(t_AA^\ast t_A^{-1}=A\). Thus the descended transpose \(\tau_A=t_A\tau\) preserves the
factorization fiber \(\mathcal F_A=\{(X,Y):YX=A\}\). Since $t_A$ is self-adjoint, $\tau_A^2=1$. More
explicitly, with \(R=S t_A^{-1}\) one has $R^t=R$, $RA=A^tR$, and
$\tau_A(X,Y)=(Y^tR,R^{-1}X^t)$. These identities follow respectively
from $St_A=t_A^tS$ and $t_AA^*t_A^{-1}=A$.
Different choices of $t_A$ differ by an element of $T_A$ and therefore induce
the same operator on $(\GL(V)\times T_A)$-invariant distributions.

Frobenius descent along the rational $T$-orbit of $A$ reduces the
fiber of $\pi_{\mathrm{ns}}$ to $\mathcal F_A$, with acting group
$\GL(V)\times T_A$ and descended transpose $\tau_A$.

If $\operatorname{rank}A=2$, then $X$ and $Y$ are isomorphisms.  Fix an
isomorphism $X_0:W\xrightarrow{\sim}V$.  Every point of $\mathcal F_A$ is
uniquely of the form
\[
    (X,Y)=\bigl(gX_0,AX_0^{-1}g^{-1}\bigr),
    \qquad g\in\GL(V).
\]
Thus $\mathcal F_A$ is one $\GL(V)$-orbit.  Its stabilizer in
$\GL(V)\times T_A$ is a diagonal copy of $T_A$, hence unimodular.  The operator $\tau_A$ preserves this rational orbit and squares to one.
Lemma~\ref{m1:lem:orbital-measures} gives its invariant-distribution statement.

Suppose next that $\operatorname{rank}A=1$.  Then $d\neq0$, and
$T_A=F^\times I_W$. With $R=S t_A^{-1}$ as above, Lemma~\ref{m1:lem:rank-one-factorization} applies and gives
transpose invariance on the rank-one factorization fiber.

It remains to treat \(\mathcal F_0=\{(X,Y)\in \Vlegs:YX=0\}\). The possible rank pairs are
\((2,0),\allowbreak(1,0),\allowbreak(0,0),\allowbreak(0,1),\allowbreak(0,2),\allowbreak(1,1)\). Let
\(U_{11}=\{(X,Y)\in\mathcal F_0: \operatorname{rank}X=\operatorname{rank}Y=1\}\). Then \(\mathcal
F_0=U_{11}\sqcup B\), with $U_{11}$ open and $B$ closed.  Lemma~\ref{m1:lem:matrix-boundary} treats
$B$.

Lemma~\ref{m1:lem:common-U11} gives the required descent. In the
trace-form model, $c(w)=\operatorname{Tr}_{K/F}(c_0w)$ for a unique
$c_0\in K^\times$.
The descended action and odd operator become
\[
 (c_0,d)\longmapsto(\alpha c_0t^{-1},\beta^{-1}td),\qquad
 (c_0,d)\longmapsto(d,c_0),
\]
with ordinary character $\eta=|\beta/\alpha|$ and odd modulus one.

Consider
\[
    \pi_{11,\mathrm{ns}}(c_0,d)=[c_0d]
    \in K^\times/F^\times.
\]
This map is invariant under the descended action.  Every fiber is one
rational orbit: if $[c_0d]=[c_0'd']$, then
$c_0'd'=r c_0d$ for some $r\in F^\times$, and the descended action above provides
$t\in K^\times$ and $\alpha,\beta\in F^\times$ carrying $(c_0,d)$ to
$(c_0',d')$.  At a representative $(1,z)$ the stabilizer is
\[
    H_z=
    \left\{
    \left(
        \begin{pmatrix}s&b\\0&s\end{pmatrix},s
    \right):
    s\in F^\times,\ b\in F
    \right\}.
\]
It is unimodular, while $\Delta_{P\times T}|_{H_z}=1$ and
$\eta|_{H_z}=1$, where $\eta=|\beta/\alpha|$. Thus the second Frobenius
compatibility character is trivial. The swap of
Lemma~\ref{m1:lem:common-U11} preserves $\eta$ and has square one.
It preserves $[c_0d]$. Lemma~\ref{m1:lem:orbital-measures} therefore fixes the positive
equivariant orbital measure.  Bernstein localization on $\pi_{11,\mathrm{ns}}$ therefore gives
\[
    \mathcal D(U_{11})^\Gamma=\mathcal D(U_{11})^{\Gamma,\tau}.
\]
An anti-invariant $\Gamma$-invariant distribution on $\mathcal F_0$ restricts to
zero on $U_{11}$ and is therefore supported on $B$, where it vanishes by
Lemma~\ref{m1:lem:matrix-boundary}.  Hence
\begin{equation}\label{m1:eq:nonsplit-zero-fiber}
    \mathcal D(\mathcal F_0)^\Gamma
    =
    \mathcal D(\mathcal F_0)^{\Gamma,\tau}.
\end{equation}

A direct two-dimensional calculation gives
\[
    \det{_F}(m_c+m_d\iota)=N_{K/F}(c)-N_{K/F}(d).
\]
Consequently,
\[
\begin{array}{c|c}
N(c)\neq N(d) & \operatorname{rank}A=2,\\
N(c)=N(d),\ (c,d)\neq(0,0) & \operatorname{rank}A=1,\\
c=d=0 & A=0.
\end{array}
\]
The three cases above show that every fiber of
$\pi_{\mathrm{ns}}$ has zero $(\widetilde \Gamma,\varepsilon)$-equivariant distribution
space.  Bernstein localization (Proposition \ref{m1:prp:localization}) gives
\[
    \mathcal D(\Vlegs)^{\widetilde \Gamma,\varepsilon}=0.
\]
This proves Theorem~\ref{m1:thm:two-leg} for the nonsplit torus.

\subsubsection{Split nonzero fibers}\label{m1:subsec:two-leg-split-nonzero}

In the split case, we take $S=I_2$ and
\(T=\{\operatorname{diag}(t_1,t_2):t_i\in F^\times\}\). We write
\[
    X=(x_1,x_2),
    \qquad
    Y=\begin{pmatrix}y_1\\y_2\end{pmatrix},
\]
where $x_i\in V$ and $y_i\in V^\vee$.  Put
\[
    u=y_1(x_1),
    \qquad
    v=y_2(x_2),
    \qquad
    \alpha=y_1(x_2),
    \qquad
    \beta=y_2(x_1),
    \qquad
    r=\alpha\beta.
\]
Then
\[
    YX=
    \begin{pmatrix}
        u&\alpha\\
        \beta&v
    \end{pmatrix},
\]
and the map
\[
    \pi_{\mathrm{sp}}:\Vlegs\longrightarrow F^3,
    \qquad
    \pi_{\mathrm{sp}}(X,Y)=(u,v,r),
\]
is $\Gamma$-invariant and fixed by transpose.

We first treat the fibers \(\pi_{\mathrm{sp}}^{-1}(u,v,r)\), \((u,v,r)\neq(0,0,0)\). Since $\det(YX)=uv-r$, the fibers divide as follows:
\[
\begin{array}{c|c}
r\neq0,\ uv-r\neq0&\text{invertible product, nonzero hyperbola}\\
r\neq0,\ uv-r=0&\text{rank-one product}\\
r=0,\ uv\neq0&\text{invertible product, complete node}\\
r=0,\ \text{exactly one of }u,v\neq0&\text{two complete nodes}\\
(u,v,r)=(0,0,0)&\text{root-component intersection and gluing}
\end{array}
\]
The last row will be treated separately.

\emph{The fiber $r\neq0$, $uv-r\neq0$.}
Then $A=YX$ is invertible, hence both $X$ and $Y$ are invertible.  For fixed
$A$, every factorization is \((X,Y)=(g,Ag^{-1})\), \(g\in\GL(V)\), so the factorization
fiber is one $\GL(V)$-orbit.  The effective split torus
acts on $(\alpha,\beta)$ by \((\alpha,\beta)\longmapsto (q\alpha,q^{-1}\beta)\), where
\(q=t_1/t_2\). Since $\alpha\beta=r\neq0$, the set of such pairs is one rational torus
orbit.  Therefore the $\pi_{\mathrm{sp}}$-fiber is a single
homogeneous $\Gamma$-orbit.  Its stabilizer is unimodular, and transpose preserves
the orbit.  The positive invariant orbital distribution is therefore fixed
by transpose, and
\[
    \mathcal D\bigl(\pi_{\mathrm{sp}}^{-1}(u,v,r)\bigr)^{\widetilde \Gamma,\varepsilon}=0.
\]

\emph{The fiber $r\neq0$, $uv-r=0$.}
Now $A=YX$ has rank one and $\alpha,\beta\neq0$.  For fixed $(u,v,r)$ the
matrices
\[
    A=\begin{pmatrix}u&\alpha\\\beta&v\end{pmatrix},
    \qquad \alpha\beta=r,
\]
form one rational $T$-orbit.  Fix one such $A_0$.  Since both off-diagonal
entries of $A_0$ are nonzero, \(T_{A_0}=F^\times I_2\). Frobenius descent along the $T$-orbit of $A_0$
reduces the fiber to
$\mathcal F_{A_0}$. If $A_0=\begin{pmatrix}u&\alpha\\\beta&v\end{pmatrix}$,
take $t_A=\operatorname{diag}(\alpha/\beta,1)$. Then
$t_AA_0^tt_A^{-1}=A_0$, and $R=t_A^{-1}$ is symmetric with
$RA_0=A_0^tR$. The descended involution is exactly $\tau_R$.
Lemma~\ref{m1:lem:rank-one-factorization} applies and gives
\[
    \mathcal D\bigl(\pi_{\mathrm{sp}}^{-1}(u,v,r)\bigr)^{\widetilde \Gamma,\varepsilon}=0.
\]

\emph{The fiber $r=0$, $uv\neq0$.}
Again $A=YX$ is invertible. In the coordinates
$(X,Y)=(g,Ag^{-1})$, the action is
$(a,t)(g,A)=(agt^{-1},tAt^{-1})$. Thus every
$\GL(V)$-invariant distribution has the form $dg\otimes D_N$,
where $dg$ is Haar measure on $\GL(V)$ and $D_N$ is a distribution
on the node $N=\{(\alpha,\beta):\alpha\beta=0\}$.
Right invariance of $dg$ identifies $T$-invariance with invariance
of $D_N$ under $(\alpha,\beta)\mapsto(q\alpha,q^{-1}\beta)$,
where $q=t_1/t_2$.
Transpose is $(g,A)\mapsto(g^{-t}A^t,A^t)$. For each fixed $A$,
inversion, transpose, and right translation preserve $dg$, so
integration in $g$ gives
$\tau_*(dg\otimes D_N)=dg\otimes(A\mapsto A^t)_*D_N$.
The latter operator exchanges $\alpha$ and $\beta$ and fixes $u,v$.
Prasad's Two-Axis Lemma therefore gives
\[
    \mathcal D\bigl(\pi_{\mathrm{sp}}^{-1}(u,v,0)\bigr)^{\widetilde \Gamma,\varepsilon}=0,
    \qquad uv\neq0.
\]

\emph{The fibers with exactly one nonzero diagonal entry.}
By symmetry assume $u\neq0$ and $v=0$.  Consider the equivariant map
\[
    p:\pi_{\mathrm{sp}}^{-1}(u,0,0)
    \longrightarrow
    \Omega_u,
    \qquad
    p(X,Y)=(x_1,y_1),
\]
where \(\Omega_u=\{(x,\ell)\in V\times V^\vee:\ell(x)=u\}\). The group $\Gamma$ acts transitively on
$\Omega_u$. We apply Frobenius descent at
$(e_1,ue_1^\vee)$. The stabilizer is
\[
    S=
    \left\{
    \left(
        \begin{pmatrix}t_1&0\\0&c\end{pmatrix},
        \operatorname{diag}(t_1,t_2)
    \right):
    t_1,t_2,c\in F^\times
    \right\}
    \simeq(F^\times)^3.
\]
Write \(x_2=xe_1+ze_2\) and \(y_2=ye_1^\vee+we_2^\vee\). The equations $v=0$ and $r=0$ become
exactly \(xy=0\) and \(zw=0\). Thus the descended fiber is \(N=N_1\times N_2\), where
\(N_1=\{(x,y):xy=0\}\) and \(N_2=\{(z,w):zw=0\}\). The two opposite-weight characters are
\(\varepsilon_1=\frac{t_1}{t_2}\) and \(\varepsilon_2=\frac{c}{t_2}\), and
\((\varepsilon_1,\varepsilon_2):S\longrightarrow(F^\times)^2\) is jointly surjective: take $t_2=1$, $t_1=q_1$, $c=q_2$ for a
prescribed pair $(q_1,q_2)$. Lemma~\ref{m1:lem:passive-products} on the
product of the two nodes gives invariance under each swap.
To track the involution, correct transpose at the chosen base by
$(\operatorname{diag}(u^{-1},1),I)$. On the fiber it acts by
\[
 (x,y,z,w)\longmapsto(u^{-1}y,ux,w,z).
\]
This is the product of the two swaps followed by the first effective
torus factor $u^{-1}$, with the second factor equal to one; both factors exist by the joint surjectivity of
$(\varepsilon_1,\varepsilon_2)$ established above. It therefore fixes every
$S$-invariant distribution. Hence
\[
    \mathcal D\bigl(\pi_{\mathrm{sp}}^{-1}(u,0,0)\bigr)^{\widetilde \Gamma,\varepsilon}=0.
\]
The case $u=0$, $v\neq0$ is obtained by interchanging the two eigenlines.

The four cases exhaust every nonzero value of $\pi_{\mathrm{sp}}$.  Therefore
\begin{equation}\label{m1:eq:split-nonzero-fibers}
    \mathcal D\bigl(\pi_{\mathrm{sp}}^{-1}(u,v,r)\bigr)^{\widetilde \Gamma,\varepsilon}=0
    \qquad
    \text{for }(u,v,r)\neq(0,0,0).
\end{equation}

\subsubsection{The split zero-product intersection}\label{m1:subsec:split-zero-intersection}

It remains to treat \(Z=\pi_{\mathrm{sp}}^{-1}(0,0,0)\). Equivalently,
\[
    Z=
    \{y_1(x_1)=y_2(x_2)=0,
      \ y_1(x_2)y_2(x_1)=0\}.
\]
Put \(Z_+=\{\beta=0\}\) and \(Z_-=\{\alpha=0\}\). Then
\[
    Z=Z_+\cup Z_-,
    \qquad
    F_0:=Z_+\cap Z_-=\{YX=0\},
\]
and transpose exchanges $Z_+$ and $Z_-$.

\begin{lemma}\label{m1:lem:split-F0}
Every $\Gamma$-invariant distribution on $F_0$ is fixed by transpose:
\[
    \mathcal D(F_0)^\Gamma=\mathcal D(F_0)^{\Gamma,\tau}.
\]
\end{lemma}

\begin{proof}
The rank pairs in $F_0$ are
\((2,0),\allowbreak(1,0),\allowbreak(0,0),\allowbreak(0,1),\allowbreak(0,2),\allowbreak(1,1)\). Let
\(U_{11}=\{(X,Y)\in F_0: \operatorname{rank}X=\operatorname{rank}Y=1\}\). Then \(F_0=U_{11}\sqcup B\),
with $U_{11}$ open and $B$ closed.  Lemma~\ref{m1:lem:matrix-boundary} treats
$B$.

By Lemma~\ref{m1:lem:common-U11}, the descended fiber has coordinates
$c=(c_1,c_2)\neq0$ and $d=(d_1,d_2)^t\neq0$, action
\eqref{m1:eq:common-U11-action}, character $\eta=|\beta/\alpha|$, and
swap $(c,d)\mapsto(d^t,c^t)$. Its odd relative modulus is one.

On the locus where \((c_1d_1,c_2d_2)\neq(0,0)\), consider
\[
    \pi_{11,\mathrm{sp}}(c,d)
    =[c_1d_1:c_2d_2]\in\mathbf P^1(F).
\]
This is invariant under the descended action and fixed by transpose.
On a fiber with both projective coordinates nonzero, if two points have
the same projective value, then for some
$\kappa\in F^\times$ one has \(c_i'd_i'=\kappa c_id_i\) for \(i=1,2\). Choose $\alpha/\beta=\kappa$ and
then choose $t_i$ so that
$c_i'=\alpha c_i/t_i$; the equations for $d_i'$ follow automatically.
Thus every fiber with both projective coordinates nonzero is a single orbit.
Its stabilizer has $\alpha=\beta$ and $t_i=\alpha$ (with $b$ arbitrary),
so it is unimodular, and both $\Delta_{P\times T}$ and $\eta$ restrict
trivially. The positive equivariant orbital distribution is fixed by the
descended involution: its action is a positive scalar of square one.

For the special fiber $[1:0]$, normalize $c_1=d_1=1$.  The remaining
variables satisfy \(c_2d_2=0\), and the residual torus acts by the standard opposite-weight action
\((c_2,d_2)\longmapsto(qc_2,q^{-1}d_2)\). The residual stabilizer has $\alpha=\beta=t_1$, with $t_2$
and $b$
free. It is unimodular, $\Delta_{P\times T}$ restricts trivially, and the
descended character is trivial. The involution fixes $c_1=d_1=1$ and
swaps $(c_2,d_2)$. Thus Prasad's
Two-Axis Lemma gives transpose invariance on this fiber.  The fiber
over $[0:1]$ is identical after interchanging the two eigenlines.

The locus omitted from $\pi_{11,\mathrm{sp}}$ consists of the two
disjoint-support orbits represented by \((c,d)=(e_1^\vee,e_2)\) and \((c,d)=(e_2^\vee,e_1)\).
For the first orbit the stabilizer is
\[
    H_z=
    \left\{
    \left(
        \begin{pmatrix}\alpha&b\\0&\beta\end{pmatrix},
        \operatorname{diag}(\alpha,\beta)
    \right)
    \right\}.
\]
The second Frobenius relative modular character is trivial, while
$\eta|_{H_z}=|\beta/\alpha|$ is nontrivial. Indeed $H_z\simeq P$ and
$\Delta_{H_z}=\Delta_{P\times T}|_{H_z}$, so there is no additional modular
factor. Hence this orbit carries no
required equivariant distribution; the transposed orbit is identical.
First apply Bernstein localization on the open domain of
$\pi_{11,\mathrm{sp}}$, using its generic and special fibers.
An odd $\eta$-equivariant distribution then has support in the omitted
closed locus. Its two disjoint-support orbits are open and closed in
that locus and both have zero required distribution space by the
character calculation above. Frobenius descent therefore gives
\[
    \mathcal D(U_{11})^\Gamma=\mathcal D(U_{11})^{\Gamma,\tau}.
\]

Now let $D\in\mathcal D(F_0)^\Gamma$ satisfy $\tau D=-D$.  Its restriction to
$U_{11}$ vanishes, so $\operatorname{supp}D\subset B$.  By
Lemma~\ref{m1:lem:matrix-boundary}, there is no $\Gamma$-invariant anti-invariant
distribution supported on $B$.  Thus $D=0$.
\end{proof}

The following consequence of the same stratification will be used in the
gluing argument.

\begin{corollary}[Boundary anomaly on $F_0$]\label{m1:cor:boundary-anomaly}
Let $Q$ be a $\tau$-anti-invariant distribution supported on $F_0$.  Suppose
that \(\operatorname{supp}(gQ-Q)\subset\{0\}\) for \(g\in \Gamma\). Then there exists $C\in\mathbf C$ such that
\begin{equation}\label{m1:eq:boundary-anomaly}
    gQ-Q=C\log\lambda(g)\,\delta_0,
\end{equation}
where, for $t=\operatorname{diag}(t_1,t_2)$, \(\lambda(a,t) =\left|\frac{\det a}{t_1t_2}\right|\). \end{corollary}

\begin{proof}
On $U_{11}$ the defect vanishes, so $Q|_{U_{11}}$ is invariant and
anti-invariant. The proof of Lemma~\ref{m1:lem:split-F0} makes this restriction
zero. Thus $Q$ is supported on $B$. Its restrictions to the two rank-two
open axes are multiples of their invariant measures. With the measures
normalized by $P_Y=\tau P_X$, anti-invariance makes their coefficients
$a,-a$. After subtracting $a(P_X-P_Y)$, the remainder is supported
on the rank-at-most-one axes and still has defect supported at zero. Hence its
restrictions to the rank-one strata are invariant, and the first part of
Lemma~\ref{m1:lem:matrix-boundary} makes them zero. The remainder is supported
at zero and is therefore a multiple of $\delta_0$, with zero defect.
Consequently
\[
 gQ-Q=-2ac\log\lambda(g)\,\delta_0,
\]
as required.
\end{proof}

\subsubsection{Root-component gluing}\label{m1:subsec:root-gluing}

Let \(D\in\mathcal D(Z)^\Gamma\) with \(\tau D=-D\). We prove that $D=0$.  Put \(U_+=Z_+\setminus
F_0\) and \(U_-=Z_-\setminus F_0\). Transpose exchanges $U_+$ and $U_-$.

\emph{Rank-mismatch orbits in $U_+$.}
On $U_+$ one has \(\beta=y_2(x_1)=0\) and \(\alpha=y_1(x_2)\neq0\). Consider
\[
    p:U_+\longrightarrow
    \Omega,
    \qquad
    p(X,Y)=(x_2,y_1),
\]
where \(\Omega=\{(x,\ell)\in V\times V^\vee:\ell(x)\neq0\}\). The group $\Gamma$ acts transitively on
$\Omega$. We apply Frobenius descent at
$(e_1,e_1^\vee)$ and write \(x_1=ae_1+de_2\) and \(y_2=be_1^\vee+ee_2^\vee\). The equations
defining $Z$ give \(a=b=0\) and \(de=0\). Hence the Frobenius fiber is \(N=\{(d,e)\in
F^2:de=0\}\). Its stabilizer is
\[
    H=
    \left\{
    \left(\operatorname{diag}(s,c),sI_2\right):
    s,c\in F^\times
    \right\},
\]
and the effective quotient acts by \(q\cdot(d,e)=(qd,q^{-1}e)\), where \(q=c/s\). The Frobenius
relative modular character is trivial.  The three strata of
$N$ correspond to
\[
\begin{array}{c|c}
(d,e) & (\operatorname{rank}X,\operatorname{rank}Y)\\ \hline
 d\neq0,\ e=0 & (2,1),\\
 d=0,\ e\neq0 & (1,2),\\
 d=e=0 & (1,1).
\end{array}
\]
The origin represents the rank-$(1,1)$ orbit in $U_+$.

Let $c_{21}^+$ and $c_{12}^+$ be the coefficients of the invariant orbital
distributions on the two punctured branches, measured relative to any
fixed positive invariant measures on the global orbits. Frobenius descent
with positive Haar measures sends positive measures to positive measures.
Thus Corollary~\ref{m1:cor:two-axis-coefficients} gives a constant $r_1>0$ with
\begin{equation}\label{m1:eq:Uplus-branch}
    c_{21}^+=r_1c_{12}^+.
\end{equation}

\emph{Comparison of the two root components.}
Consider the $\Gamma$-invariant open subset \(U_X=\{(X,Y)\in Z:\det X\neq0\}\). The projection
$U_X\to\GL_2(F)$, $(X,Y)\mapsto X$, has homogeneous base.
Frobenius descent allows us to fix $X=I_2$.  The stabilizer is the diagonal
copy of $T$, and the fiber is
\[
    N_X=
    \left\{
    \begin{pmatrix}0&\alpha\\\beta&0\end{pmatrix}:
    \alpha\beta=0
    \right\}.
\]
The effective torus acts by \((\alpha,\beta)\longmapsto(q\alpha,q^{-1}\beta)\). The two punctured
branches correspond to the rank-$(2,1)$ orbits in $U_+$
and $U_-$, respectively. The same positivity argument applied to this
chart gives a constant $r_2>0$ with
\begin{equation}\label{m1:eq:Xchart-branch}
 c_{21}^+=r_2c_{21}^-.
\end{equation}
Transpose sends the rank-$(2,1)$ orbit in $U_-$ to the
rank-$(1,2)$ orbit in $U_+$. Its pushforward preserves positivity, so anti-invariance gives
some $r_3>0$ with
\begin{equation}\label{m1:eq:transpose-branch}
 c_{21}^-=-r_3c_{12}^+.
\end{equation}
Combining \eqref{m1:eq:Uplus-branch}, \eqref{m1:eq:Xchart-branch}, and
\eqref{m1:eq:transpose-branch} gives $(r_1+r_2r_3)c_{12}^+=0$. Since all three constants are positive, this
implies, also using transpose for the fourth orbit,
\[
    c_{21}^+=c_{12}^+=c_{21}^-=c_{12}^-=0.
\]
Thus all rank-mismatch contributions on $Z\setminus F_0$ vanish.

\emph{The remaining rank-one coefficient.}
Define the closed linear root components \(S_+=\{x_1=0,\ y_2=0\}\) and \(S_-=\{x_2=0,\
y_1=0\}\). They are exchanged by transpose and satisfy \(S_+\cap S_-=\{0\}\). The open set $S_+\cap U_+$
is precisely the rank-$(1,1)$ orbit in $U_+$,
and $S_-\cap U_-$ is its transpose.

On $S_+$ write a point as $(u,\varphi)\in V\times V^\vee$ and put
\[
    q_+(u,\varphi)=\varphi(u),
    \qquad
    \rho(a,t)=\left|\frac{t_1}{t_2}\right|,
    \quad
    t=\operatorname{diag}(t_1,t_2).
\]
For $\Re(s)\gg0$, define
\[
    Z_+(s)(f)=
    \int_{S_+}f(u,\varphi)|\varphi(u)|^s\,du\,d\varphi,
\]
and define \(Z_-(s)=\tau Z_+(s)\). We view these distributions on $Z$ through the closed embeddings
$S_\pm\hookrightarrow Z$.

The determinant family in Lemma~\ref{m1:lem:matrix-boundary} also supplies
the analytic properties needed here. Indeed, the linear isomorphism
\begin{equation}\label{m1:eq:determinant-pairing-isomorphism}
 J(u,\varphi)=\begin{pmatrix}u_1&u_2\\-\varphi_2&\varphi_1\end{pmatrix},
 \qquad \det J(u,\varphi)=\varphi(u),
\end{equation}
preserves normalized product Haar measure. It therefore identifies the
test-function integrals for $Z_+(s)$ with determinant zeta integrals.
Consequently $Z_+(s)$ has a simple pole at $s=-2$, with residue
$R=c\delta_0$ for the same $c\ne0$, and is regular there away from the
origin. Since $Z_-(s)=\tau Z_+(s)$, its residue is also $R$.
This identifies the meromorphic families as distributions; the
transformation character for the action on $S_+$ is computed next.


The action of $g=(a,t)\in \Gamma$ on $S_+$ is
\[
    (u,\varphi)
    \longmapsto
    \bigl(t_2^{-1}au,\,t_1\varphi a^{-1}\bigr).
\]
Thus $q_+$ is multiplied by $t_1/t_2$ and the linear map has
absolute determinant $\rho(g)^2$. Under the pushforward convention
the product Haar measure therefore transforms by $\rho(g)^{-2}$.  Consequently
\[
    gZ_+(s)=\rho(g)^{-(s+2)}Z_+(s).
\]
Write
\[
    Z_+(s)=\frac{R}{s+2}+P_++O(s+2).
\]
Then
\begin{equation}\label{m1:eq:Pplus-defect}
    gP_+-P_+=-\log\rho(g)\,R.
\end{equation}
Choose $P_-=\tau P_+$.  Since conjugation by $\tau$ sends $(a,t)$ to $(a^{-t},t^{-1})$ and therefore
inverts $\rho$,
\[
    gP_--P_-=+\log\rho(g)\,R.
\]
Hence
\begin{equation}\label{m1:eq:Ppm-defect}
    g(P_+-P_-)-(P_+-P_-)
    =-2\log\rho(g)\,R.
\end{equation}

Let \(O_{11}^+=S_+\cap U_+\) and \(O_{11}^-=S_-\cap U_-\). On $O_{11}^+$ one has $q_+\neq0$, so
$Z_+(s)$ is regular there at $s=-2$.
Thus \(\nu_+:=P_+|_{O_{11}^+}\) is a nonzero invariant orbital distribution.  Put \(\nu_-:=\tau\nu_+\);
then $\nu_-=P_-|_{O_{11}^-}$.

Since all rank-mismatch contributions vanish, there exists
$a\in\mathbf C$ such that \(D|_{U_+}=a\nu_+\) and \(D|_{U_-}=-a\nu_-\). Define
\(Q=D-a(P_+-P_-)\). Then \(\operatorname{supp}Q\subset F_0\). Since $D$ is $\Gamma$-invariant,
\eqref{m1:eq:Ppm-defect} gives
\begin{equation}\label{m1:eq:Q-rho-defect}
    gQ-Q=2a\log\rho(g)\,R
    =2ac\log\rho(g)\,\delta_0.
\end{equation}
On the other hand, Corollary~\ref{m1:cor:boundary-anomaly} gives
\[
    gQ-Q=C\log\lambda(g)\,\delta_0
\]
for some $C\in\mathbf C$.  Thus
\[
    2ac\log\rho(g)=C\log\lambda(g)
    \qquad(g\in \Gamma).
\]
The characters $\rho$ and $\lambda$ are independent.  Indeed, let $\varpi$
be a uniformizer of $F$ and take
\[
    g=\bigl(I_2,\operatorname{diag}(\varpi,\varpi^{-1})\bigr).
\]
Then
\[
    \lambda(g)=1,
    \qquad
    \rho(g)=|\varpi|^2\neq1.
\]
Since $c\neq0$, it follows that $a=0$.

Therefore $D|_{Z\setminus F_0}=0$, so $\operatorname{supp}D\subset F_0$.
By Lemma~\ref{m1:lem:split-F0}, every $\Gamma$-invariant distribution supported on
$F_0$ is fixed by transpose.  Since $\tau D=-D$, we conclude that $D=0$.
Hence
\begin{equation}\label{m1:eq:split-zero-fiber}
    \mathcal D(Z)^{\widetilde \Gamma,\varepsilon}=0.
\end{equation}

By \eqref{m1:eq:split-nonzero-fibers} and \eqref{m1:eq:split-zero-fiber}, every
fiber of $\pi_{\mathrm{sp}}$ has zero $(\widetilde \Gamma,\varepsilon)$-equivariant
distribution space.  The map $\pi_{\mathrm{sp}}$ is continuous and
$\widetilde \Gamma$ preserves every fiber.  Bernstein localization
(Proposition~\ref{m1:prp:localization}) therefore gives
\[
    \mathcal D(\Vlegs)^{\widetilde \Gamma,\varepsilon}=0.
\]
Together with the nonsplit calculation, this proves
Theorem~\ref{m1:thm:two-leg}.

\subsection{The linear identity descendant}\label{m1:sec:identity}
\subsubsection{The representation and nonzero quadratic fibers}
\label{m1:subsec:identity-fibers}
Let $\Gamma=\GL_2(F)\times\GL_2(F)$ act on
$\mathcal N=M_2(F)^2\oplus\mathfrak{gl}_2(F)$ by
\begin{equation}\label{m1:eq:identity-action}
 (a,b)(X,Y,Z)=(aXb^{-1},bYa^{-1},bZb^{-1}),\qquad
 \tau(X,Y,Z)=(Y^t,X^t,Z^t).
\end{equation}
The involution normalizes $\Gamma$ by $(a,b)\mapsto(a^{-t},b^{-t})$;
write $\widetilde\Gamma=\Gamma\rtimes\langle\tau\rangle$ and use the
sign character $\varepsilon$.

\begin{theorem}[Identity Descendant]\label{m1:thm:identity-descendant}
Every $\Gamma$-invariant distribution on $\mathcal N$ is fixed by $\tau$; equivalently,
$\mathcal D(\mathcal N)^{\widetilde\Gamma,\varepsilon}=0$.
\end{theorem}
The proof first localizes the sign space to a quadratic null cone. The
same localization applies after partial Fourier transform, giving the
second support condition needed for odd-dimensional homogeneity.

Since $F$ has characteristic zero, write
$Z=zI+Z_0$, where $z=\tr Z/2$ and $Z_0\in W:=\mathfrak{sl}_2(F)$.
Thus $\mathcal N=M_2(F)^2\times F\times W$. The scalar $z$ is fixed
by both actions, whereas the legs and traceless variable retain their
actions. On $W$ use
\[
 B(U,V)=\tr(UV),\qquad q(U)=B(U,U).
\]
The symmetric trace pairing is nondegenerate: for
$U=\left(\begin{smallmatrix}a&b\\c&-a\end{smallmatrix}\right)$ it pairs
as $2aa'+bc'+cb'$. Both $B$ and $q$ are invariant under conjugation
and transpose.

\begin{proposition}\label{m1:prop:identity-nonzero-fibers}
For every $z\in F$ and $c\in F^\times$, the fiber
\[
 \mathcal N_{z,c}=M_2(F)^2\times\{z\}\times\mathcal O_c,
 \qquad\mathcal O_c=\{Z_0\in W:q(Z_0)=c\},
\]
has zero sign-equivariant distribution space.
\end{proposition}
\begin{proof}
Every $Z_0\in\mathcal O_c$ has characteristic polynomial $T^2-c/2$.
For $c\neq0$ this polynomial is separable, so $Z_0$ is regular
semisimple. Rational canonical form shows that $\mathcal O_c$ is one
rational $\GL_2(F)$-orbit. Its centralizer
$T_{Z_0}=F[Z_0]^\times$ is the split or nonsplit maximal torus according
to this polynomial. Frobenius descent to a chosen $Z_0$ leaves the
two-leg representation
\[
 \GL_2(F)\times T_{Z_0}\curvearrowright M_2(F)^2,
 \qquad(a,t)(X,Y)=(aXt^{-1},tYa^{-1}).
\]
Choose the rational symmetric form $S$ from
Section~\ref{m1:subsec:two-leg-statement}: the standard form in a split
eigenbasis or the trace pairing on $F[Z_0]$ in the nonsplit case.
Then $SZ_0=Z_0^tS$ and $St=t^tS$ for all $t\in T_{Z_0}$. The
correcting element $(1,S^{-1})$ sends $Z_0^t$ to $Z_0$, and the induced
operator on the fiber is
\[
 \tau_{Z_0}(X,Y)=(Y^tS,S^{-1}X^t)=\tau_S(X,Y).
\]
By \eqref{m1:eq:two-leg-adjoint}, it has square one and normalizes the
two-leg action by $(a,t)\mapsto(a^{-t},t^{-1})$. Both the ambient ordinary group and
$\GL_2(F)\times T_{Z_0}$ are unimodular. Their finite extensions are
unimodular, so the relative modulus is trivial on the odd component
as well. The descended character is the sign character, and
Theorem~\ref{m1:thm:two-leg} gives the assertion.
\end{proof}

\begin{corollary}\label{m1:cor:identity-support}
Every $T\in\mathcal D(\mathcal N)^{\widetilde\Gamma,\varepsilon}$
satisfies
\begin{equation}\label{m1:eq:identity-first-support}
 \supp T\subset M_2(F)^2\times F\times Z(B),
 \qquad Z(B)=\{U\in W:B(U,U)=0\}.
\end{equation}
\end{corollary}
\begin{proof}
On the invariant open set $q(Z_0)\neq0$, the continuous invariant map
$(X,Y,z,Z_0)\mapsto(z,q(Z_0))$ has target $F\times F^\times$.
Proposition~\ref{m1:prop:identity-nonzero-fibers} treats every
fiber, so Bernstein localization makes the restriction of $T$ to this
open set zero. The closed complement gives the stated support.
\end{proof}
For traceless $2\times2$ matrices, $q(U)=-2\det U$. Thus $Z(B)$ is
the nilpotent cone in $W$.

\subsubsection{Fourier equivariance and homogeneity}
\label{m1:subsec:identity-Fourier}
Let $dV$ be the self-dual Haar measure on $W$ for $B$ and the fixed
additive character $\psi$. For
$f\in C_c^\infty(M_2(F)^2\times F\times W)$, set
\begin{equation}\label{m1:eq:partial-Fourier}
 (\mathcal F_W f)(X,Y,z,U)
 =\int_W f(X,Y,z,V)\psi(B(U,V))\,dV,
 \qquad (\mathcal F_WT)(f)=T(\mathcal F_Wf).
\end{equation}

\begin{lemma}\label{m1:lem:identity-Fourier}
Partial Fourier transform preserves
$\mathcal D(\mathcal N)^{\widetilde\Gamma,\varepsilon}$. In particular
every distribution in this space also satisfies
\begin{equation}\label{m1:eq:identity-second-support}
 \supp(\mathcal F_WT)\subset M_2(F)^2\times F\times Z(B).
\end{equation}
\end{lemma}
\begin{proof}
Conjugation and transpose preserve the pairing on $W$:
\[
 \tr(gUg^{-1}gVg^{-1})=\tr(UV),\qquad
 \tr(U^tV^t)=\tr(VU)=\tr(UV).
\]
An isometry of a nondegenerate symmetric form has determinant of
absolute value one and preserves the self-dual measure. Changing
variables in \eqref{m1:eq:partial-Fourier} therefore proves that these
operators commute with $\mathcal F_W$ on functions and distributions.
The action on the remaining variables is independent of the $W$
integration. Thus $\mathcal F_W$ commutes with the extended group action.
Corollary~\ref{m1:cor:identity-support}, applied to $\mathcal F_WT$, now gives
\eqref{m1:eq:identity-second-support}.
\end{proof}
\begin{proof}[Proof of Theorem~\ref{m1:thm:identity-descendant}]
We apply the common homogeneity theorem, Theorem~\ref{m2:tools:homogeneity}.

Here $W=\mathfrak{sl}_2(F)$ with the nondegenerate trace form and
$M=M_2(F)^2\times F$, an affine analytic manifold. Equations~\eqref{m1:eq:identity-first-support} and
\eqref{m1:eq:identity-second-support} give both support hypotheses, with
the Fourier convention \eqref{m1:eq:partial-Fourier}. Since $\dim_F W=3$,
the theorem gives $T=0$.
\end{proof}

\subsection{Invariance of distributions on the linear normal spaces}\label{m1:sec:descendants}
\label{m1:subsec:dense-descendants}\label{m1:subsec:rankzero-descendants}
\label{m1:subsec:rankone-descendants}\label{m1:subsec:exhaustive-descendants}

\begin{proposition}\label{m1:prop:all-descendants}
\label{m1:thm:dense-descendants}\label{m1:thm:rank00-descendants}
\label{m1:thm:decomposable-descendant}\label{m1:thm:stable-system-descendant}
For every closed rational orbit $E\in\cX(F)$, let
$N_E=T_E\cX/T_E(LE)$, and let $\widetilde L_E$ be the extended
stabilizer with its restricted sign character. Then
\[
 \mathcal D(N_E)^{\widetilde L_E,\varepsilon_E}=0.
\]
\end{proposition}
\begin{proof}
Theorem~\ref{m1:thm:closed-orbits} and
Section~\ref{m1:subsec:normal-interface} give the nine normal
representations with their corrected odd operators. Throughout,
Lemma~\ref{m1:lem:passive-products} allows us to retain the summands
fixed by the extended stabilizer as parameters.

\emph{Dense orbits.}
For the absolutely irreducible type the odd operator is the identity.
The split toral type is an opposite-weight plane with rationally
surjective character $t_1/t_2$, so
Proposition~\ref{m1:prop:opposite-weight-plane} applies.
For the nonsplit type, use the rational root module
\eqref{m1:eq:dense-nonsplit-quotient} with action
\eqref{m1:eq:dense-nonsplit-action}. Hilbert~90 identifies the effective
action of $K^\times/F^\times$ with all of $K^1$. Thus each nonzero
norm fiber is a single $K^\times$-orbit with stabilizer $F^\times$;
both groups are unimodular. The odd map is $d\mapsto-\bar d$. Since $-1\in K^1$,
it acts on invariant distributions as conjugation does, and
Lemma~\ref{m2:elem:norm} applies to the root module.
For the scalar pair, Lemma~\ref{m1:lem:matrix-pair-transpose} applies
to the normal space $M_2(F)^2$.

\emph{Rank-zero orbits.}
For regular $D$, split or nonsplit,
Section~\ref{m1:subsec:rankzero-normal} identifies the nontrivial part
of the normal representation and its odd operator with the adjoint
Two-Leg representation; the two
centralizer directions are fixed. Theorem~\ref{m1:thm:two-leg}
therefore applies. For $D=dI$, $d\ne0$, the action
\eqref{m1:eq:rankzero-scalar} is the
Identity Descendant of Theorem~\ref{m1:thm:identity-descendant}.

\emph{Rank-one orbits.}
For a decomposable rank-one orbit, the normal-space formulas
\eqref{m1:eq:decomposable-normal}--\eqref{m1:eq:decomposable-transpose}
show that the characters $(\alpha,\beta)$ are jointly rationally surjective.
The invariant map on the root coordinates is
\begin{equation}\label{m1:eq:decomposable-fiber-map}
 \pi(x,y,u,v)=(xy,uv),\qquad
 \pi^{-1}(c_1,c_2)=\{xy=c_1\}\times\{uv=c_2\}.
\end{equation}
Each factor is a hyperbola or the two-axis node. Lemma~\ref{m1:lem:passive-products} and Bernstein
localization exclude sign-equivariant distributions on the normal space. For a stable-system rank-one orbit, the stabilizer character $\eta$
is rationally surjective
by \eqref{m1:eq:stable-stabilizer};
\eqref{m1:eq:stable-ext-table} and
Lemma~\ref{m1:lem:stable-parameter-basis} give one opposite-weight
plane with an ordinary swap and four fixed directions.
Proposition~\ref{m1:prop:opposite-weight-plane} therefore applies
throughout $b_1=1-ad+sb_0$, $\det H\ne0$, including nilpotent and
repeated-root values of $D$.

These cases exhaust the closed rational orbits. Conjugation transports
the normal representations and extended stabilizers from the chosen
representatives to every point of each rational orbit.
\end{proof}

\subsection{Proof of the linear multiplicity-one theorem}\label{m1:sec:completion}
\begin{proof}[Proof of Theorem~\ref{m1:thm:main}]
The residual variety $\cX\simeq\GL_4$ is smooth affine and the action
is algebraic. Corollary~\ref{m1:cor:transpose-stable} gives rational
transpose stability of its closed orbits, and
Proposition~\ref{m1:prop:all-descendants} excludes sign-equivariant
distributions on their normal spaces. Corollary~\ref{m1:cor:HC-invariance} therefore yields
\[
 \mathcal D(\cX)^{\widetilde L,\varepsilon}=0.
\]
Thus every residual invariant distribution is transpose-invariant.

The unimodular, involution-compatible Frobenius isomorphism
\eqref{m1:eq:residual-Frobenius} carries this conclusion back to all
$H_{\mathrm L}\times H_{\mathrm L}$-invariant distributions on
$G_{\mathrm L}=\GL_4(F)\times\GL_2(F)$. Proposition~\ref{m1:prop:GK} gives, for every
irreducible admissible smooth complex representation $\pi$,
\begin{equation}\label{m1:eq:final-product-bound}
 m_{H_{\mathrm L}}(\pi)m_{H_{\mathrm L}}(\pi^\vee)\leq1,
 \qquad m_{H_{\mathrm L}}(\pi)=\dim\Hom_{H_{\mathrm L}}(\pi,\Co).
\end{equation}

To identify the two multiplicities, let
$\iota_{G_{\mathrm L}}(g,z)=(g^{-t},z^{-t})$. The contragredient theorem for general
linear groups~\cite[Theorem~7.3]{BZ76}, applied to each factor, gives
$\pi^\vee\simeq\pi\circ\iota_{G_{\mathrm L}}$. Moreover
\[
 \iota_{G_{\mathrm L}}(\diag(a,b),b)=(\diag(a^{-t},b^{-t}),b^{-t}),
\]
so $\iota_{G_{\mathrm L}}$ preserves the embedded $H_{\mathrm L}$ and restricts to an automorphism
of it. The $H_{\mathrm L}$-invariant linear forms on $\pi\circ\iota_{G_{\mathrm L}}$ are exactly
those on $\pi$, and hence
\begin{equation}\label{m1:eq:final-multiplicity-equality}
 m_{H_{\mathrm L}}(\pi^\vee)=m_{H_{\mathrm L}}(\pi).
\end{equation}
Combining \eqref{m1:eq:final-product-bound} and
\eqref{m1:eq:final-multiplicity-equality} gives $m_{H_{\mathrm L}}(\pi)^2\leq1$, and
therefore $m_{H_{\mathrm L}}(\pi)\leq1$.
\end{proof}
